% Options for packages loaded elsewhere
\PassOptionsToPackage{unicode}{hyperref}
\PassOptionsToPackage{hyphens}{url}
\documentclass[
]{article}
\usepackage{xcolor}
\usepackage[margin=1in]{geometry}
\usepackage{amsmath,amssymb}
\setcounter{secnumdepth}{-\maxdimen} % remove section numbering
\usepackage{iftex}
\ifPDFTeX
  \usepackage[T1]{fontenc}
  \usepackage[utf8]{inputenc}
  \usepackage{textcomp} % provide euro and other symbols
\else % if luatex or xetex
  \usepackage{unicode-math} % this also loads fontspec
  \defaultfontfeatures{Scale=MatchLowercase}
  \defaultfontfeatures[\rmfamily]{Ligatures=TeX,Scale=1}
\fi
\usepackage{lmodern}
\ifPDFTeX\else
  % xetex/luatex font selection
\fi
% Use upquote if available, for straight quotes in verbatim environments
\IfFileExists{upquote.sty}{\usepackage{upquote}}{}
\IfFileExists{microtype.sty}{% use microtype if available
  \usepackage[]{microtype}
  \UseMicrotypeSet[protrusion]{basicmath} % disable protrusion for tt fonts
}{}
\makeatletter
\@ifundefined{KOMAClassName}{% if non-KOMA class
  \IfFileExists{parskip.sty}{%
    \usepackage{parskip}
  }{% else
    \setlength{\parindent}{0pt}
    \setlength{\parskip}{6pt plus 2pt minus 1pt}}
}{% if KOMA class
  \KOMAoptions{parskip=half}}
\makeatother
\usepackage{color}
\usepackage{fancyvrb}
\newcommand{\VerbBar}{|}
\newcommand{\VERB}{\Verb[commandchars=\\\{\}]}
\DefineVerbatimEnvironment{Highlighting}{Verbatim}{commandchars=\\\{\}}
% Add ',fontsize=\small' for more characters per line
\usepackage{framed}
\definecolor{shadecolor}{RGB}{248,248,248}
\newenvironment{Shaded}{\begin{snugshade}}{\end{snugshade}}
\newcommand{\AlertTok}[1]{\textcolor[rgb]{0.94,0.16,0.16}{#1}}
\newcommand{\AnnotationTok}[1]{\textcolor[rgb]{0.56,0.35,0.01}{\textbf{\textit{#1}}}}
\newcommand{\AttributeTok}[1]{\textcolor[rgb]{0.13,0.29,0.53}{#1}}
\newcommand{\BaseNTok}[1]{\textcolor[rgb]{0.00,0.00,0.81}{#1}}
\newcommand{\BuiltInTok}[1]{#1}
\newcommand{\CharTok}[1]{\textcolor[rgb]{0.31,0.60,0.02}{#1}}
\newcommand{\CommentTok}[1]{\textcolor[rgb]{0.56,0.35,0.01}{\textit{#1}}}
\newcommand{\CommentVarTok}[1]{\textcolor[rgb]{0.56,0.35,0.01}{\textbf{\textit{#1}}}}
\newcommand{\ConstantTok}[1]{\textcolor[rgb]{0.56,0.35,0.01}{#1}}
\newcommand{\ControlFlowTok}[1]{\textcolor[rgb]{0.13,0.29,0.53}{\textbf{#1}}}
\newcommand{\DataTypeTok}[1]{\textcolor[rgb]{0.13,0.29,0.53}{#1}}
\newcommand{\DecValTok}[1]{\textcolor[rgb]{0.00,0.00,0.81}{#1}}
\newcommand{\DocumentationTok}[1]{\textcolor[rgb]{0.56,0.35,0.01}{\textbf{\textit{#1}}}}
\newcommand{\ErrorTok}[1]{\textcolor[rgb]{0.64,0.00,0.00}{\textbf{#1}}}
\newcommand{\ExtensionTok}[1]{#1}
\newcommand{\FloatTok}[1]{\textcolor[rgb]{0.00,0.00,0.81}{#1}}
\newcommand{\FunctionTok}[1]{\textcolor[rgb]{0.13,0.29,0.53}{\textbf{#1}}}
\newcommand{\ImportTok}[1]{#1}
\newcommand{\InformationTok}[1]{\textcolor[rgb]{0.56,0.35,0.01}{\textbf{\textit{#1}}}}
\newcommand{\KeywordTok}[1]{\textcolor[rgb]{0.13,0.29,0.53}{\textbf{#1}}}
\newcommand{\NormalTok}[1]{#1}
\newcommand{\OperatorTok}[1]{\textcolor[rgb]{0.81,0.36,0.00}{\textbf{#1}}}
\newcommand{\OtherTok}[1]{\textcolor[rgb]{0.56,0.35,0.01}{#1}}
\newcommand{\PreprocessorTok}[1]{\textcolor[rgb]{0.56,0.35,0.01}{\textit{#1}}}
\newcommand{\RegionMarkerTok}[1]{#1}
\newcommand{\SpecialCharTok}[1]{\textcolor[rgb]{0.81,0.36,0.00}{\textbf{#1}}}
\newcommand{\SpecialStringTok}[1]{\textcolor[rgb]{0.31,0.60,0.02}{#1}}
\newcommand{\StringTok}[1]{\textcolor[rgb]{0.31,0.60,0.02}{#1}}
\newcommand{\VariableTok}[1]{\textcolor[rgb]{0.00,0.00,0.00}{#1}}
\newcommand{\VerbatimStringTok}[1]{\textcolor[rgb]{0.31,0.60,0.02}{#1}}
\newcommand{\WarningTok}[1]{\textcolor[rgb]{0.56,0.35,0.01}{\textbf{\textit{#1}}}}
\usepackage{graphicx}
\makeatletter
\newsavebox\pandoc@box
\newcommand*\pandocbounded[1]{% scales image to fit in text height/width
  \sbox\pandoc@box{#1}%
  \Gscale@div\@tempa{\textheight}{\dimexpr\ht\pandoc@box+\dp\pandoc@box\relax}%
  \Gscale@div\@tempb{\linewidth}{\wd\pandoc@box}%
  \ifdim\@tempb\p@<\@tempa\p@\let\@tempa\@tempb\fi% select the smaller of both
  \ifdim\@tempa\p@<\p@\scalebox{\@tempa}{\usebox\pandoc@box}%
  \else\usebox{\pandoc@box}%
  \fi%
}
% Set default figure placement to htbp
\def\fps@figure{htbp}
\makeatother
\setlength{\emergencystretch}{3em} % prevent overfull lines
\providecommand{\tightlist}{%
  \setlength{\itemsep}{0pt}\setlength{\parskip}{0pt}}
\usepackage{bookmark}
\IfFileExists{xurl.sty}{\usepackage{xurl}}{} % add URL line breaks if available
\urlstyle{same}
\hypersetup{
  pdftitle={MCAR{,} MAR{,} MNAR},
  hidelinks,
  pdfcreator={LaTeX via pandoc}}

\title{MCAR, MAR, MNAR}
\author{}
\date{\vspace{-2.5em}}

\begin{document}
\maketitle

\textbf{Inference lab} using the five-step template: Hypothesis →
Visualise → Assumptions → Conduct → Conclude.

\subsection{Learning objectives}\label{learning-objectives}

\begin{itemize}
\tightlist
\item
  Define MCAR, MAR, and MNAR and give an example of each.
\item
  Simulate the three mechanisms and measure the bias in a complete-case
  analysis.
\item
  Decide when a missingness assumption is defensible from the data alone
  and when it is not.
\end{itemize}

\subsection{Prerequisites}\label{prerequisites}

Regression basics.

\subsection{Background}\label{background}

The language of missingness is Rubin's taxonomy. MCAR (missing
completely at random) means the probability of being missing does not
depend on anything --- observed or unobserved. MAR (missing at random)
means the probability depends only on observed variables. MNAR (missing
not at random) means the probability depends on the missing value
itself, even after conditioning on everything observed. Each is a
property of the process, not the data, and you can rarely decide between
MAR and MNAR from the data alone.

Complete-case analysis drops any row with a missing value. It is
unbiased under MCAR, usually biased under MAR, and always suspect under
MNAR. Multiple imputation handles MAR correctly if the imputation model
is compatible with the analysis model, but cannot rescue MNAR without
explicit modelling of the missingness mechanism.

Auxiliary variables --- things that correlate with both the outcome and
the missingness but are not in the analysis model --- are the currency
of MAR. If you have them, put them in the imputation model. If you do
not, your MAR claim is on thinner ice than it appears.

\subsection{Setup}\label{setup}

\begin{Shaded}
\begin{Highlighting}[]
\FunctionTok{library}\NormalTok{(tidyverse)}
\FunctionTok{set.seed}\NormalTok{(}\DecValTok{42}\NormalTok{)}
\FunctionTok{theme\_set}\NormalTok{(}\FunctionTok{theme\_minimal}\NormalTok{(}\AttributeTok{base\_size =} \DecValTok{12}\NormalTok{))}
\end{Highlighting}
\end{Shaded}

\subsection{1. Hypothesis}\label{hypothesis}

We estimate the mean of Y from each of three datasets with the same true
mean, where Y has been made missing by MCAR, MAR, and MNAR processes.

\subsection{2. Visualise}\label{visualise}

\begin{Shaded}
\begin{Highlighting}[]
\NormalTok{df }\OtherTok{\textless{}{-}} \FunctionTok{tibble}\NormalTok{(}
  \AttributeTok{x =} \FunctionTok{rnorm}\NormalTok{(n, }\DecValTok{50}\NormalTok{, }\DecValTok{10}\NormalTok{),}
  \AttributeTok{y =} \DecValTok{2} \SpecialCharTok{+} \FloatTok{0.5} \SpecialCharTok{*}\NormalTok{ x }\SpecialCharTok{+} \FunctionTok{rnorm}\NormalTok{(n, }\DecValTok{0}\NormalTok{, }\DecValTok{5}\NormalTok{)}
\NormalTok{)}

\NormalTok{mcar }\OtherTok{\textless{}{-}}\NormalTok{ df }\SpecialCharTok{|\textgreater{}} \FunctionTok{mutate}\NormalTok{(}\AttributeTok{y\_obs =} \FunctionTok{if\_else}\NormalTok{(}\FunctionTok{runif}\NormalTok{(n) }\SpecialCharTok{\textless{}} \FloatTok{0.30}\NormalTok{, }\ConstantTok{NA\_real\_}\NormalTok{, y),}
                     \AttributeTok{mech  =} \StringTok{"MCAR"}\NormalTok{)}
\NormalTok{mar  }\OtherTok{\textless{}{-}}\NormalTok{ df }\SpecialCharTok{|\textgreater{}} \FunctionTok{mutate}\NormalTok{(}\AttributeTok{p =} \FunctionTok{plogis}\NormalTok{(}\SpecialCharTok{{-}}\DecValTok{2} \SpecialCharTok{+} \FloatTok{0.05} \SpecialCharTok{*}\NormalTok{ (x }\SpecialCharTok{{-}} \DecValTok{50}\NormalTok{)),}
                     \AttributeTok{y\_obs =} \FunctionTok{if\_else}\NormalTok{(}\FunctionTok{runif}\NormalTok{(n) }\SpecialCharTok{\textless{}}\NormalTok{ p, }\ConstantTok{NA\_real\_}\NormalTok{, y),}
                     \AttributeTok{mech  =} \StringTok{"MAR"}\NormalTok{)}
\NormalTok{mnar }\OtherTok{\textless{}{-}}\NormalTok{ df }\SpecialCharTok{|\textgreater{}} \FunctionTok{mutate}\NormalTok{(}\AttributeTok{p =} \FunctionTok{plogis}\NormalTok{(}\SpecialCharTok{{-}}\DecValTok{2} \SpecialCharTok{+} \FloatTok{0.05} \SpecialCharTok{*}\NormalTok{ (y }\SpecialCharTok{{-}} \FunctionTok{mean}\NormalTok{(y))),}
                     \AttributeTok{y\_obs =} \FunctionTok{if\_else}\NormalTok{(}\FunctionTok{runif}\NormalTok{(n) }\SpecialCharTok{\textless{}}\NormalTok{ p, }\ConstantTok{NA\_real\_}\NormalTok{, y),}
                     \AttributeTok{mech  =} \StringTok{"MNAR"}\NormalTok{)}

\NormalTok{all }\OtherTok{\textless{}{-}} \FunctionTok{bind\_rows}\NormalTok{(}
\NormalTok{  mcar }\SpecialCharTok{|\textgreater{}} \FunctionTok{select}\NormalTok{(mech, x, y, y\_obs),}
\NormalTok{  mar  }\SpecialCharTok{|\textgreater{}} \FunctionTok{select}\NormalTok{(mech, x, y, y\_obs),}
\NormalTok{  mnar }\SpecialCharTok{|\textgreater{}} \FunctionTok{select}\NormalTok{(mech, x, y, y\_obs)}
\NormalTok{)}

\NormalTok{all }\SpecialCharTok{|\textgreater{}}
  \FunctionTok{mutate}\NormalTok{(}\AttributeTok{missing =} \FunctionTok{is.na}\NormalTok{(y\_obs)) }\SpecialCharTok{|\textgreater{}}
  \FunctionTok{ggplot}\NormalTok{(}\FunctionTok{aes}\NormalTok{(x, y, }\AttributeTok{colour =}\NormalTok{ missing)) }\SpecialCharTok{+}
  \FunctionTok{geom\_point}\NormalTok{(}\AttributeTok{alpha =} \FloatTok{0.5}\NormalTok{) }\SpecialCharTok{+}
  \FunctionTok{facet\_wrap}\NormalTok{(}\SpecialCharTok{\textasciitilde{}}\NormalTok{ mech) }\SpecialCharTok{+}
  \FunctionTok{labs}\NormalTok{(}\AttributeTok{colour =} \StringTok{"Y missing?"}\NormalTok{)}
\end{Highlighting}
\end{Shaded}

\subsection{3. Assumptions}\label{assumptions}

Complete-case analysis is the bluntest possible method. It is justified
if missingness is MCAR. We will compute the mean of Y on observed rows
for each mechanism and compare with the true mean.

\subsection{4. Conduct}\label{conduct}

\begin{Shaded}
\begin{Highlighting}[]
  \FunctionTok{group\_by}\NormalTok{(mech) }\SpecialCharTok{|\textgreater{}}
  \FunctionTok{summarise}\NormalTok{(}
    \AttributeTok{truth =} \FunctionTok{mean}\NormalTok{(y),}
    \AttributeTok{cc    =} \FunctionTok{mean}\NormalTok{(y\_obs, }\AttributeTok{na.rm =} \ConstantTok{TRUE}\NormalTok{),}
    \AttributeTok{bias  =}\NormalTok{ cc }\SpecialCharTok{{-}}\NormalTok{ truth,}
    \AttributeTok{pct\_missing =} \FunctionTok{mean}\NormalTok{(}\FunctionTok{is.na}\NormalTok{(y\_obs)) }\SpecialCharTok{*} \DecValTok{100}
\NormalTok{  )}
\NormalTok{results}
\end{Highlighting}
\end{Shaded}

\begin{Shaded}
\begin{Highlighting}[]
\NormalTok{results }\SpecialCharTok{|\textgreater{}}
  \FunctionTok{ggplot}\NormalTok{(}\FunctionTok{aes}\NormalTok{(mech, bias)) }\SpecialCharTok{+}
  \FunctionTok{geom\_col}\NormalTok{(}\AttributeTok{alpha =} \FloatTok{0.7}\NormalTok{, }\AttributeTok{fill =} \StringTok{"steelblue"}\NormalTok{) }\SpecialCharTok{+}
  \FunctionTok{geom\_hline}\NormalTok{(}\AttributeTok{yintercept =} \DecValTok{0}\NormalTok{, }\AttributeTok{linetype =} \StringTok{"dashed"}\NormalTok{) }\SpecialCharTok{+}
  \FunctionTok{labs}\NormalTok{(}\AttributeTok{x =} \ConstantTok{NULL}\NormalTok{, }\AttributeTok{y =} \StringTok{"Bias of complete{-}case mean"}\NormalTok{)}
\end{Highlighting}
\end{Shaded}

\subsection{5. Concluding statement}\label{concluding-statement}

\begin{quote}
Complete-case estimation of the mean of Y was unbiased under MCAR (bias
≈ \texttt{round(results\$bias{[}results\$mech\ ==\ "MCAR"{]},\ 2)}) and
meaningfully biased under MAR
(\texttt{round(results\$bias{[}results\$mech\ ==\ "MAR"{]},\ 2)}) and
MNAR (\texttt{round(results\$bias{[}results\$mech\ ==\ "MNAR"{]},\ 2)}),
despite the same true mean and similar missingness proportions.
\end{quote}

\subsection{Common pitfalls}\label{common-pitfalls}

\begin{itemize}
\tightlist
\item
  Using Little's test as evidence \emph{for} MCAR.
\item
  Imputing with a model that omits the analysis outcome.
\item
  Reporting a sample size inflated by a complete-case filter that hides
  40\% dropout.
\item
  Using ``the missingness was not significantly related to X'' as proof
  of MAR.
\end{itemize}

\subsection{Further reading}\label{further-reading}

\begin{itemize}
\tightlist
\item
  Little RJA, Rubin DB. \emph{Statistical Analysis with Missing Data}.
\item
  van Buuren S (2018), \emph{Flexible Imputation of Missing Data}.
\item
  Sterne JA et al.~(2009), \emph{Multiple imputation for missing data in
  epidemiological and clinical research}.
\end{itemize}

\subsection{Session info}\label{session-info}

\begin{Shaded}
\begin{Highlighting}[]

\end{Highlighting}
\end{Shaded}

\subsection{See also --- chapter index}\label{see-also-chapter-index}

\begin{itemize}
\tightlist
\item
  \href{../index.Rmd}{Methods in this chapter}
\item
  \href{../../../ch00-find-your-method.Rmd}{Find your method (Chapter
  0)}
\end{itemize}

\end{document}
